% Generated draft: requires mathematical and source review.
\documentclass[11pt]{article}
\usepackage[margin=25mm]{geometry}
\usepackage{fontspec}
\setmainfont{Latin Modern Roman}
\usepackage{amsmath,amssymb,mathtools,unicode-math}
\setmathfont{Latin Modern Math}
\usepackage{enumitem,needspace,xurl,hyperref}
\setlength{\parindent}{0pt}
\setlength{\parskip}{5pt}
\setlength{\emergencystretch}{3em}
\newcommand{\N}{\mathbb{N}}
\newcommand{\Z}{\mathbb{Z}}
\newcommand{\Q}{\mathbb{Q}}
\newcommand{\R}{\mathbb{R}}
\newcommand{\C}{\mathbb{C}}
\newcommand{\F}{\mathbb{F}}
\newcommand{\A}{\mathbb{A}}
\newcommand{\PP}{\mathbb{P}}
\newcommand{\E}{\mathbb{E}}
\newcommand{\eps}{\varepsilon}
\newcommand{\dd}{\,\mathrm{d}}
\begin{document}
\section*{UnsolvedMath ID: 30000496}
\textbf{Source problem code:} catalogue entry
\subsection{Definitions and mathematical statement}
The source asks to classify the finite groups \(G\) admitting Beauville structures, equivalently the groups occurring in its intended Beauville-surface construction. The target is a classification of groups up to isomorphism, not a classification of all structures on each group or all resulting surfaces.

In the usual unmixed complex-geometry convention, a Beauville structure consists of two generating pairs \((x_1,y_1)\), \((x_2,y_2)\) in \(G\). Put \(z_i=(x_iy_i)^{-1}\), and define
\[
 \Sigma_i=\bigcup_{g\in G}\bigcup_{m\geq1}
 \{gx_i^mg^{-1},\ gy_i^mg^{-1},\ gz_i^mg^{-1}\}.
\]
The conditions are \(\langle x_i,y_i\rangle=G\) for \(i=1,2\), the disjointness condition
\[
 \Sigma_1\cap\Sigma_2=\{1\},
\]
and the hyperbolic triangle condition
\[
 \frac1{\operatorname{ord}(x_i)}+
 \frac1{\operatorname{ord}(y_i)}+
 \frac1{\operatorname{ord}(z_i)}<1\quad(i=1,2)
\]
for the genus-at-least-two curve interpretation. This spells out a standard unmixed realization problem: determine exactly the finite groups for which such data exist, with completeness of the classification and no restriction to a preselected family of finite groups.

Geometrically, the usual unmixed construction gives a free diagonal action of \(G\) on \(C_1\times C_2\), where each \(C_i\) is a smooth complex curve of genus at least two and \(C_i/G\) is a projective line with three branch values, yielding \((C_1\times C_2)/G\). The group-theoretic disjointness condition expresses freeness of the diagonal action. This conventional description must not be assumed to cover a mixed construction, in which the full group can exchange the two factors. The supplied wording does not distinguish unmixed and mixed Beauville structures, so the displayed conditions are explicitly conditional on the standard unmixed interpretation; a mixed classification needs its different subgroup/action data.

The literature assessment checked 2026-08-21 says no mathematical question was recoverable beyond an incomplete-answer remark. That does not match the explicit finite-group classification question in both current statement fields. The verification field says the classification request was preserved and its construction reference made self-contained. This mismatch is recorded rather than used to infer a current classification result or open status. The open label remains dated metadata.
\subsection{Short English statement}
Classify the finite groups that support the specified Beauville construction. The standard unmixed formulation uses two generating pairs whose conjugate cyclic-power sets meet only in the identity, with hyperbolic triangle data. The source does not distinguish this from mixed constructions, and its older literature note does not match the explicit classification question.
\subsection{Sources}
\begin{itemize}
\item Local UnsolvedMath registry, permanent ID 30000496, source code catalogue entry, ``Finite Groups Admitting Beauville Structures,'' including the original surface-construction question, cleaned group-classification request and verification metadata.
\item \emph{Pro-\(p\) Extensions of Global Fields and pro-\(p\) Groups (2007)}, Oberwolfach Reports, pp.\ 1463--1536, DOI 10.4171/owr/2006/25: \url{https://doi.org/10.4171/owr/2006/25}. This is the sole supplied literature reference, labelled ``original OWR report.'' The citation uses 2007, while the DOI and proposed-year metadata use 2006.
\item Literature triage checked 2026-08-21; local record updated 2026-09-16. No external source has been visited or independently checked here.
\end{itemize}
\end{document}
